\Page{14}{\textenglish{03 / CAMPAIGNS}}{ישראל, מסטילס לסרט}{42}{48}

\vspace{5mm}

{\fontsize{13.6}{21}\selectfont

\StoryColumns{\textbf{רנואר · ארבעה פריימים ויומן אחד.} אנה זק ויעל שלביה היו עסוקות בפסטיגל כשרנואר ביקשה להפיק איתן קמפיין. דניאלה קרול סיפרה ל״פורטפוליו״ שהצוות יצר ארבע תמונות סטילס, שתיים לכל אחת, בסביבות אדריכליות נקיות. כלי ההתאמה של \textenglish{Astria} שימש להכנת הדימויים, והצוות של \textenglish{FutureCreators} הוביל את הקריאייטיב וההפקה. קרול תיארה את העבודה כתרגום של השפה הקיימת של רנואר לחומר חזותי חדש. לפי התיעוד שלה, זק ושלביה נתנו אישור ותובנות, בלי להשתתף ביום צילום לקמפיין הזה. הפוסטים והדימויים הופיעו בערוצי המותג והפרזנטוריות. התנאים המסחריים והעלויות לא פורסמו. \SourceRef{9}

במקרה כזה חשוב לשמור על מה שהקהל אמור לזהות. קיר או זווית אפשר לשנות כדי להוציא את הלוק מהיומיום. הפנים ופרטי הבגד הם חלק מן ההבטחה של הפריים. במקרה הזה, הכלי היה דרך להמשיך לעבוד למרות יומן צפוף, והאתגר היה להגיע לתמונות שנראות שייכות דווקא לרנואר.}{\textbf{\textenglish{H\&O} · המצלמה זזה בתוך סצנה מחוללת.} בספטמבר 2026 השיקה הרשת בישראל את \textenglish{Real Life Brands} לקולקציית חורף 2027. סרט הקמפיין מציג רגעים רגילים, כמו יציאה עם הכלב או הדרך לפגישה, כסצנות אופנה. לפי תיאור ההפקה של \textenglish{Kanibal Media}, הפריימים נוצרו בכלי \textenglish{AI} בלי ימי צילום עם דוגמנים מצולמים. הצוות בנה גם סיבובי מצלמה תלת־ממדיים סביב סצנות קפואות והשלים אותן באנימציה ובעריכה. בקרדיטים מופיעים במאי, יוצר \textenglish{AI}, מפיק וצוותי קריאייטיב ופוסט. הצילומים הווירטואליים עדיין דרשו החלטות של בימוי, תנועה וקצב. \SourceRef{9}}

\par}

\vspace{5mm}

\FeatureIntro{PaperPink}{שני המקרים מראים למה חשוב לדעת מה בדיוק הופק. ברנואר צריך להחזיק את הזהות ואת הלוק בארבע תמונות. ב־\textenglish{H\&O} צריך גם לשמר את הסצנה כשהזווית מתחלפת. ארבעה סטילס וסרט מחולל יכולים לדרוש צוותים שונים, גם כשהכותרת בשני המקרים היא ״קמפיין \textenglish{AI}״.}
